%=============================================================================!
% 18.1  DATA, MODEL AND LOSS                                                  !
%=============================================================================!
\section{Data, model and loss}
\label{sec:lo-loss}

A spring is stretched by 1, 2, \dots, \SI{10}{\centi\metre}, and at each
stretch the force is read on a dial that errs by about
\SI{0.05}{\newton} (\cref{fig:lo-bond}a, readings drawn by computer for a
spring of \SI{25}{\newton\per\metre}). Hooke's law (\cref{sec:ne-spring})
says the force is $kx$ for some stiffness $k$, and the task is to choose
$k$. Three things make up any such task. The \emph{data} are the pairs
of stretch and force, $(x_i, F_i)$ for $i = 1$ to $n$; the \emph{model}
is the rule $F = kx$ that turns a stretch into a predicted force, with
$k$ its one adjustable number, its \emph{weight}; and the \emph{loss}
measures how far the model misses the data. The usual loss is the sum of
squared misses,
\begin{equation}
   \mathcal{L}(k) = \sum_{i=1}^n\bigl(F_i - kx_i\bigr)^2 ,
   \label{eq:lo-loss}
\end{equation}
which is zero only if every reading lies on the line, is never negative,
and counts a miss twice as large four times as heavily. Learning, here,
is finding the $k$ that makes $\mathcal{L}$ least.

$\mathcal{L}$ is a parabola in $k$ (\cref{fig:lo-bond}b): expanding the square,
$\mathcal{L} = \sum_iF_i^2 - 2k\sum_ix_iF_i + k^2\sum_ix_i^2$, a quadratic with a
positive coefficient of $k^2$. Its lowest point is where its slope
vanishes:
\begin{align*}
   \frac{\dd\mathcal{L}}{\dd k} &= -2\sum_ix_i\bigl(F_i - kx_i\bigr) = 0
      && \text{differentiating each square by the chain rule}\\
   \sum_ix_iF_i &= k\sum_ix_i^2
      && \text{dividing by $-2$ and moving the $k$ terms}\\
   \hat k &= \frac{\sum_ix_iF_i}{\sum_ix_i^2}
      && \text{solving, the hat marking the fitted value.}
\end{align*}
For the readings of the figure $\hat k = \SI{25.72}{\newton\per\metre}$.
Its error follows as in the least-squares line of \cref{sec:cv-drift}:
$\hat k$ is a weighted sum of the readings, $\sum_iw_iF_i$ with $w_i =
x_i/\sum_jx_j^2$, so if the readings err independently by
$\sigma_{\mathrm m}$, its variance is $\sigma_{\mathrm
m}^2\sum_iw_i^2 = \sigma_{\mathrm m}^2/\sum_ix_i^2$. With
$\sigma_{\mathrm m}$ estimated from the misses about the fitted line,
\SI{0.056}{\newton}, the stiffness is $(25.72 \pm
0.28)$\,\si{\newton\per\metre}, 2.5 errors from the
\SI{25}{\newton\per\metre} the readings were drawn from. Ten readings
fix the stiffness to about 1\%, and with an error estimated from only ten
readings a miss this large comes in 3\% of such fits, by Student's
factor (\cref{sec:cv-mean}).

\subsection*{The stiffness of a bond}

The same fit gives the stiffness of a chemical bond from a molecular
dynamics run. The A--B molecule of Chapter~16, two atoms of nitrogen's
mass held by a Morse well of depth $D = \SI{1}{\electronvolt}$ and width
$a = \SI{2}{\per\angstrom}$, has the stiffness $2Da^2 =
\SI{8}{\electronvolt\per\angstrom\squared}$ at the bottom of its well
(\cref{sec:os-anharmonic}). Run for \SI{20}{\pico\second} under Langevin
dynamics, each \SI{10}{\femto\second} gives a reading of the stretch $x
= r - r_0$ and the force along the bond, 2001 readings in all
(\cref{fig:lo-bond}c). At \SI{100}{\kelvin} the stretch spreads by
\SI{0.0345}{\angstrom}, close to the $\sqrt{\kB T/k} =
\SI{0.0328}{\angstrom}$ that equipartition gives a spring of that
stiffness (\cref{sec:sm-equipartition}); at \SI{1000}{\kelvin} by
\SI{0.127}{\angstrom}, against \num{0.104}, reaching
\SI{0.72}{\angstrom} on the soft outer side of the well.

A straight line through these readings needs two weights, an intercept
$\theta_0$ and a slope, $F \approx \theta_0 - kx$, since the bond's mean length need
not sit at $r_0$. Two weights are found as one was, with a little more
bookkeeping, which the matrix form of the next toolbox does once for any
number of weights. The line's stiffness is \SI{7.66}{\electronvolt\per
\angstrom\squared} at \SI{100}{\kelvin}, 7.07 at 300 and 5.26 at
\SI{1000}{\kelvin}, below the 8 of the well's bottom and further below
the hotter the run. The fit is sound: the Morse well's own
stiffness changes along it, as $2Da^2(2\mathrm{e}^{-2ax} -
\mathrm{e}^{-ax}) \approx 8(1 - 3ax)$ for small $x$, by about a fifth
over the range the bond explores even at \SI{100}{\kelvin}, and a line
fitted over a range gives the stiffness averaged over that range, which
the softer outer side pulls down. A model answers the question its data
pose: readings from a hot run teach the hot bond. Adding the powers $x^2$,
$x^3$ and $x^4$ to the model lets the fitted curve bend with the force,
and its slope at $x = 0$ then gives 8.126, 8.007 and
\SI{8.000}{\electronvolt\per\angstrom\squared} at \SI{100}{\kelvin}, and
\num{8.084} at \SI{1000}{\kelvin} with four powers. Choosing what goes
into a model is the subject of \cref{sec:lo-basis}.

\begin{toolbox}[Least squares in matrix form]
A model whose prediction is a sum of weights times known functions,
$\hat y(x) = \sum_{j=1}^p\theta_j\phi_j(x)$, is linear in its weights. At the
$n$ readings $x_i$ the predictions form the vector $X\boldsymbol{\theta}$, with the
\emph{design matrix} $X_{ij} = \phi_j(x_i)$, one row per reading and one
column per function; for the line $\theta_0 - kx$, the columns are 1 and
$x_i$, and $\boldsymbol{\theta} = (\theta_0, -k)$. The loss is the squared length of the vector of
misses,
\begin{equation*}
   \mathcal{L}(\boldsymbol{\theta}) = (X\boldsymbol{\theta} - \vb{y})^{\mathsf T}(X\boldsymbol{\theta} - \vb{y})
   = \sum_{i=1}^n\Bigl(\sum_{l=1}^pX_{il}\theta_l - y_i\Bigr)^2 .
\end{equation*}
Its rate of change with one weight $\theta_j$, by the chain rule on each
square, is
\begin{equation*}
   \frac{\partial\mathcal{L}}{\partial\theta_j}
   = \sum_i2\Bigl(\sum_lX_{il}\theta_l - y_i\Bigr)X_{ij}
   = 2\bigl[X^{\mathsf T}(X\boldsymbol{\theta} - \vb{y})\bigr]_j ,
\end{equation*}
since $\sum_iX_{ij}m_i$ is the $j$th entry of $X^{\mathsf T}\vb{m}$ for
any vector of misses $\vb{m}$. The $p$ rates together are the
\emph{gradient of the loss}, $\grad\mathcal{L} = 2X^{\mathsf T}(X\boldsymbol{\theta} -
\vb{y})$, a vector pointing the way $\mathcal{L}$ rises fastest in the space of
weights (\cref{sec:en-gradient}). At the least loss every rate vanishes,
which gives the \emph{normal equations}
\begin{equation}
   X^{\mathsf T}X\,\boldsymbol{\theta} = X^{\mathsf T}\vb{y} ,
   \label{eq:lo-normal}
\end{equation}
$p$ linear equations for $p$ weights, solved as in \cref{sec:ro-remove}.
For one weight and the single column $x_i$ they are
$\bigl(\sum_ix_i^2\bigr)k = \sum_ix_iy_i$, the result above.
\end{toolbox}

\begin{figure}[htb]
\centering
\includegraphics{ch18_learning/bond}
\caption{One weight learned by least squares. (a) A spring's force read
on a dial with an error of \SI{0.05}{\newton} at ten stretches (readings
drawn by computer for \SI{25}{\newton\per\metre}), and the line through
the origin of least squares. (b) The sum of squared misses against the
stiffness tried; the fit is its lowest point (dashed). (c) The A--B
molecule of Chapter~16 under Langevin dynamics: the force along the bond
against its stretch, every \SI{50}{\femto\second} of \SI{20}{\pico
\second}, at 100 and \SI{1000}{\kelvin}, with the Morse force (grey
dashed) and the straight lines of least squares.}
\label{fig:lo-bond}
\end{figure}

\begin{keyidea}
Learning a model from data is the lowering of a loss, here the sum of
squared misses. For a model linear in its weights the least loss solves
the normal equations $X^{\mathsf T}X\boldsymbol{\theta} = X^{\mathsf T}\vb{y}$ at
once. A fit answers the question its data pose: a line through the forces
of a hot run gives the stiffness averaged over the stretches that run
explores.
\end{keyidea}

\notebook{18}{fit a spring and a bond, and change the temperature of the
bond's run}

\begin{selfcheck}
Ten readings of a spring's force at stretches $x_i$ all lie exactly on a
line through the origin. What is $\mathcal{L}(\hat k)$, and what does the formula
for the error of $\hat k$ then give?
\end{selfcheck}
